% Imported from LaTeX_and_pdfs/AGAT_dump.tex; see provenance.json.
\part{Graphs, Finite Graphs}


\section{Graphs, Finite Graphs, Trees }

Serre (1980) \cite{agat:Serr1980}  

cf. Chapter I. Trees and Amalgams, Section 1 Amalgams, Subsection 1.1 Direct limits of Serre (1980) \cite{agat:Serr1980}  


Let $(G_i)_{i\in I}$, family of groups.    

$\forall \, $ pair $(i,j)$, let $F_{ij} = $ set of homomorphisms of $G_i$ into $G_j$

Want: group $G= \varinjlim G_i$ and 
\[
\lbrace f_i | f_i : G_i \to G \rbrace \text{s.t. } f_j \circ f = f_i \quad \, \forall \, f \in F_{ij}
\]
group $G$ and family $\lbrace f_i\rbrace$ universal in that  

(*) if $H$ group, if $\lbrace h_i | h_i :G_i \to H ; h_j \circ f = h_i \qquad \, \forall \, f \in F_{ij} \rbrace$, \\
then $\exists \, ! h: G\to H$ s.t. $h_i = h\circ f_i$  \\
i.e. $\text{Hom}(G,H) \simeq \varprojlim \text{Hom}(G_i,H)$, the inverse limit being taken relative to $F_{ij}$.  \\
i.e. $G$ direct limit of $G_i$ relative to the $F_{ij}$.  

EY : 20170918 this is my rewrite/reinterpretation:

Let $(G_i)_{i\in I}$, $\forall \, (i,j) \in I^2$, let $F_{ij} = \lbrace f \equiv f_{ij} | f:G_i \to G_j , f \text{ homomorphism of $G_i$ into $G_j$ } \rbrace$.  

Given group $G = \varinjlim G_i$ (for fixed $i$), $\lbrace f_i | f_i : G_i \to G  |  f_j \circ f = f_i \quad \, \forall \, f \in F_{ij} \rbrace$, i.e.  

\begin{tikzpicture}
\matrix(m)[matrix of math nodes, row sep=4em, column sep=8em]
{
G_i    &  G   \\
G_j   &  \\};
\path[->]
(m-1-1) edge node[auto]{$f_i$} (m-1-2)
edge node[auto]{$f_{ij} \equiv f$} (m-2-1) 
(m-2-1) edge node[right]{$ f_j$ } (m-1-2);
\end{tikzpicture} 

Then $G$,  $\lbrace f_i | f_i : G_i \to G  |  f_j \circ f = f_i \quad \, \forall \, f \in F_{ij} \rbrace$ \textbf{universal} \\
\phantom{Then } if $\forall \, $ group $H$, $\forall \, \lbrace h_i |  h_i : G_i \to H  | h_j \circ f = h_i \quad \, \forall \, f \in F_{ij} \rbrace$, \\
\phantom{Then if } then $\exists \, ! \  h : G\to H$, s.t. $h_i = h \circ f_i$ i.e. 
\begin{tikzpicture}
\matrix(m)[matrix of math nodes, row sep=4em, column sep=8em]
{
H & G_i    &  G   \\
& G_j   &  \\};
\path[->]
(m-1-2) edge node[auto]{$f_i$} (m-1-3)
edge node[auto]{$f_{ij} \equiv f$} (m-2-2) 
edge node [auto]{$h_i$} (m-1-1)
(m-2-2) edge node[right]{$ f_j$ } (m-1-3)
edge node [auto]{$h_j$} (m-1-1) 
(m-1-3) edge[bend right = 30] node [auto]{ $\exists \, ! h$ } (m-1-1)  
;
\end{tikzpicture} 






\begin{proposition}
	$\exists \, !$ pair $G$, family $(f_i)_{i\in I}$, i.e. (pair consisting of $G, (f_i)_{i\in I}$, unique up to unique isomorphism.  
\end{proposition}
\begin{proof}
Define $G$ by generators and relations.  \\
Take generating family to be disjoint union of those for $G_i$.  \\
relations - $xyz^{-1}$ where $x,y,z \in G_i$, $z=xy \in G_i$ \\
\phantom{relations - } $xy^{-1}$ where $x\in G_i$, $y \in G_j$, $y=f(x)$ for at least $f\in F_{ij}$.  

Thus, existence of $G,\lbrace f_i\rbrace$.  

$G$ represents functor $H\mapsto \varprojlim \text{Hom}{(G_i,H)}$.  

Thus, uniqueness (also from universal property).  
\end{proof}

e.g. groups $A,G_1,G_2$, homomorphisms $\begin{aligned} & \quad \\ 
& f_1 : A \to G_1 \\
& f_2 : A \to G_2 
\end{aligned}$.  

$G$ obtained by amalgamating $A$ in $G_1,G_2$ by $f_1,f_2 \equiv G_1 *_A G_2$.  \\
1 can have $G=\lbrace 1 \rbrace$, even though $f_1,f_2$ non-trivial.  

\emph{Application}: (Van Kampen Thm.)

Let topological space $X$ be covered by open $U_1,U_2$.   \\
Suppose $U_1,U_2, U_{12}=U_1\bigcap U_2$ arcwise connected.  

Let basept. $x\in U_{12}$.  

Then $\pi_1(X;x)$ obtained by taking 3 groups 
\[
\pi_1(U_1;x), \pi_1(U_2;x), \pi_1(U_{12};x)
\]
and amalagamating them according to homomorphism
\[
\begin{aligned}
& \pi_1(U_{12};x) \to \pi_1(U_1;x)  \\
& \pi_1(U_{12};x) \to \pi_1(U_2;x)
\end{aligned}
\]

\exercisehead{1}  
Let homomorphisms $\begin{aligned} & \quad \\
	& f_1 : A \to G_1 \\ 
	& f_2: A \to G_2 \end{aligned}
$
amalgam $G=G_1 *_A G_2$.  

Define subgroups $A^n,G_1^n, G_2^n$, of $A,G_1,G_2$ recursively by 
\[
\begin{aligned}
& A^1 = \lbrace 1 \rbrace \\
& G_1^1 = \lbrace 1 \rbrace \\ 
& G_2^1 = \lbrace 1 \rbrace \\ 
\end{aligned}
\]

$A^n = $ subgroup of $A$ generated by $f_1^{-1}(G_1^{n-1})$ and $f_2^{-1}(G_2^{n-1})$  
\[
G_1^n = \text{subgroup of $G_i$ generated by $f_i(A^n)$ }
\]

Let $A^{\infty}, G_i^{\infty}$ be unions of $A^n, G_i^n$ resp.  

Show that $f_i$ defines injection $A/A^{\infty} \to G_i/G_i^{\infty}$.  

So the amalgamation is $G \simeq G_1/G_1^{\infty} *_{A/A^{\infty}} G_2/G_2^{\infty}$.  

Take the first induction case (for intuition about the solution).  

\[
\begin{aligned}
	& A^2 = \langle f_1^{-1}( G_1^1), f_2^{-1}(G_2^1) \rangle = \langle f_1^{-1}(\lbrace 1 \rbrace), f_2^{-1}(\lbrace 1 \rbrace ) \rangle \\
	& G_i^2 = f_i(A^2)
\end{aligned}
\]
Let $f_i(a) = f_i(b) \in G_i/G_i^{\infty}$; $a,b\in A/A^{\infty}$.  

Then since $f_i(a),f_i(b) \in G_i/G_i^{\infty}$, $f_i(a),f_i(b) \in \lbrace gG_i^{\infty} | g\in G_i \rbrace$ (quotient is defined to be the set of all left cosets of $G_i^{\infty}$, which has to be a normal subgroup for $G_i/G_i^{\infty}$ to be a quotient group).  

Since $a,b \in A/A^{\infty}$, suppose we take $a,b\in A$.  

And suppose we take 
\[
\begin{aligned}
& 	f_i(a) = f_i(a)G_i^{\infty} = f_i(a) f_i(A^{n_a} ) = f_i(aA^{n_a}) \\ 
& 	f_i(b) = f_i(b)G_i^{\infty} = f_i(b) f_i(A^{n_b} ) = f_i(bA^{n_b})  
\end{aligned}
\]
Taking $f_i^{-1}$ (recall for group homomorphisms, they map inverse of element of 1st. group to inverse of image of this element).  

$aA^{n_a}=bA^{n_b}\in A/A^{\infty}$ (This is okay as we've "quotiented out $A^{\infty}$; so indeed, they're equal)


cf. Subsection 1.2 Structure of amalgams of Serre (1980) \cite{agat:Serr1980} 

Suppose given group $A$, family of groups $(G_i)_{i\in I}$, and, $\forall \, i\in I$, injective homomorphism $A\to G_i$.  

$*_A G_i \equiv $ direct limit (cf. no. 1.1) of family $(A,G_i)$ with respect to these homomorphisms, call it \emph{sum} (in category theory sense, i.e. product) of $G_i$ with $A$ amalgamated.  

e.g. $A=\lbrace 1 \rbrace$, \\
$*G_i \equiv $ free product of $G_i$.  

\subsubsection{reduced word}  

$\forall \, i \in I$, choose set $S_i$ of right coset representations of $G_i$ modulo $A$, 

assume $1 \in S_i$, 

$(a,s)\mapsto as$ is bijection of $A\times S_i$ onto $G_i$,  \\
$A\times (S_i-\lbrace 1 \rbrace) \to G_i-A$ (onto)

Let $\mathbf{i} = (i_1\dots i_n)$, $n\geq 0$, $i_j \in I$, s.t. 
\begin{equation}
i_m \neq i_{m+1} \text{ for } 1 \leq m \leq n-1
\end{equation}
cf. (T) of Serre (1980) \cite{agat:Serr1980}.  

So \emph{reduced word} $m$ is defined as 
\[
m = (a;s_1\dots s_n)
\]
where $a\in A, s_1\in S_{i_1} \dots s_n \in S_{i_n}$, and $s-j \neq 1\, \forall \, j$.  

$f\equiv $ canonical homomorphism of $A$ into group $G= *_A G_i$ \\ 
$f_i \equiv $ canonical homomorphism of $G_i$ into group $G= *_A G_i$

EY : 20170611 (Further explanations, basic examples, from me):  

Given $A, \lbrace G_i\rbrace_{i\in I}$, injective (group) homomorphisms $\lbrace f_i: A \to G_i\rbrace_i$.  

$G_i \backslash f_i(A) = \lbrace f_i(A)g | g\in G_i\rbrace$.  

Right coset representation of $f_i(A)g\mapsto g$.  

e.g. $A,G_1,G_2$, $\begin{aligned} & \quad \\
	& f_1:A \to G_1 \\
		& f_2 : A\to G_2 \end{aligned}$.  
		
	\[
	\begin{aligned}
	& G_1\backslash f_1(A) = \lbrace f_1(A)g| g\in G_1\rbrace \\
	& G_2\backslash f_2(A) = \lbrace f_2(A)g | g\in G_2 \rbrace
	\end{aligned}
	\]

$\mathbf{i} = (i_1\dots i_n)$, $i_j\in I$, $i_m\neq i_{m+1}$ for $1\leq m \leq n-1$.  

Consider $(1212\dots 12)$  

$m=(a;f_1 g_2 f_3 g_4 \dots f_{2n-1}, g_{2n})$ where $f$'s $\in S_1 \subset G_1$, $g$'s $\in S_2 \subset G_2$.  

and so 
\begin{definition}[reduced word]
	\textbf{reduced word} of type $\mathbf{i}$, $m$, \begin{equation}
	m=(a;s_1\dots s_n)
	\end{equation}
	where $a\in A, s_1 \in S_{i_1}, \dots s_n \in S_{i_n}$, $s_j\neq 1$ \, $\forall \, j$, \\
	\phantom{where } $\mathbf{i} = (i_1\dots i_n)$, $i_j \in I$, s.t. $i_m \neq i_{m+1}$ for $1\leq m \leq n-1$, \\
	with $S_i = \lbrace g | g\in f_i(A)g \in f_i(A) G_i\rbrace$  
\end{definition}




\begin{theorem}[1 of Serre (1980) \cite{agat:Serr1980} ]
	$\forall \, g \in G$, $\exists \, $ sequence $\mathbf{i}$ s.t. $i_m \neq i_{m+1}$ for $1\leq m \leq n-1$ and 
	
	reduced word 
	\[
	m = (a;s_1\dots s_n) 
	\]
	of type $\mathbf{i}$ s.t. 
	\[
	g = f(a)f_{i_1}(s_1) \dots f_{i_n}(s_n)
	\]
\end{theorem}

Furthermore, $\mathbf{i}$ and $m$ unique.  

\emph{Remark}.  Thm. 1 implies $f;f_i$ injective.  

Then identify $A$ and $G_i$ with images $f(A), f_i(G_i)$ in $G$, and reduced decomposition (*) of $g\in G$ 
\[
g = as_1\dots s_n, \quad \, a\in A, \, s_1 \in S_{i_1} - \lbrace 1 \rbrace \dots s_n \in S_{i_n} - \lbrace 1 \rbrace
\]
Likewise, $G_i \bigcap G_j = A$ if $i\neq j$.  

In particular, $S_i - \lbrace 1 \rbrace$ pairwise disjoint in $G$.  

\begin{proof}
Let $X_i \equiv $set of reduced words of type $\mathbf{i}$, $X = \coprod X_i$.  

Make $G$ act on $X$.  

In view of universal property of $G$, sufficient to make $\forall \, i, G_i$ act, 

check action induced on $A$ doesn't depend on $i$  

Suppose then that $i\in I$, and let $Y_i = $ set of reduced words of form $(1;s_1 \dots s_n)$, with $i_1\neq i$.  

EY : 20170611

Recall that 
\[
S_i = \lbrace g| g\in f_i(A) g \in f_i(A) G_i \rbrace
\]
\[
\begin{aligned}
	& A \times S_i \to G_i \text{ onto } \\ 
	& A\times (S_i - \lbrace 1 \rbrace) \to G_i - A \text{ onto } \\ 
	& (a,s) \mapsto as \text{ bijection }
\end{aligned}
\]

Let $Y_i = $ set of reduced words of form $(1; s_1 \dots s_n) = \lbrace (1;s_1 \dots s_n) | 1\in A; s_1 \in S_{i_1}\dots s_n \in S_{i_n} ; \mathbf{i} = (i_1\dots i_n), \, i_j \in I \text{ s.t. } i_m \neq i_{m+1} \text{ for } 1\leq m \leq n -1 \rbrace$.  

\[
\begin{gathered}
\begin{gathered}
A \times Y_i \to X = \coprod_i X_i \\
(a,(1; s_1 \dots s_n)) \mapsto (a; s_1 \dots s_n)
\end{gathered} \\
\begin{gathered}
A\times \lbrace S_i - \lbrace 1 \rbrace ) \times Y_i \to X \\
((a,s), (1;s_1 \dots s_n)) \mapsto (a;s,s_1\dots s_n)
\end{gathered}
\end{gathered}
\]
and remember that $X_i = $ set of reduced words of type $\mathbf{i}$.  

It's clear that this yields a bijection $A\times Y_i \bigcup A\times (S_i - \lbrace 1 \rbrace) \times Y_i \to X$.  

Let $x\in X$.  Then $x\in X_{\mathbf{i}}$ for some $\mathbf{i}$.  So $x$ is a reduced word of type $\mathbf{i}$: $x = (a;s_1\dots s_n)$.  Then clearly $x = (a;s_1\dots s_n) \mapsto (a,(1;s_1\dots s_n)) \in A\times Y_i$.  So $\forall \, x \in X$, $\exists (a, (1; s_1 \dots s_n) \in A \times Y_i)$. So we've proven, for part of bijection, the mapping going one way. \\

Identify $A \bigcup A \times (S_i - \lbrace 1 \rbrace )$ with $G_i$, with \\
\phantom{Identify} $a \in A$, \, $A \times (S_i - \lbrace 1 \rbrace) \ni (a, s_i)$. \\

$G_i$ act on $G_i \times Y_i$ :  \\
\[
\begin{aligned}
G_i \times (G_i \times Y_i) & \mapsto G_i \times Y_i \\
g' \cdot (g, y) & = (g' g , y)
\end{aligned}
\]

(group $G_i, Y_i \equiv $ set of reduced words) \\

Transfer this to an action of $G_i$ on $X$ by means of $\theta_i$: \\
\[
\begin{aligned}
G_i \times X & \mapsto X \\
g' \cdot (a, s_1, \dots , s_n) & = (g'a, s_1 , \dots s_n) 
\end{aligned}
\]
$X \equiv $ set of reduced words of type $\mathbf{i}$. \\

restriction to $A \Longrightarrow a' \cdot (a; s_1 , \dots s_n ) = (a'a ; s_1, \dots s_n)$, which doesn't depend on $i$.

We've now defined action of $G$ on $X$.

\end{proof}


cf. pp. 13, Sec. 2. Trees, 2.1 Graphs of Serre (1980) \cite{agat:Serr1980}

\begin{definition}[1. of Serre (1980) \cite{agat:Serr1980}]
	\textbf{graph} $\Gamma = (X,Y, Y\to X\times X, Y\to Y)$, where $\begin{aligned} & \quad \\
		& \text{ set } X = \text{ vert } \Gamma \\ 
		& \text{ set } Y = \text{ edge } \Gamma \end{aligned}$  

\[
\begin{aligned}
& Y \to X\times X \\ 
& y\mapsto (o(y), t(y)) \\
& Y\to Y \\
& y\mapsto \overline{y}
\end{aligned}
\]
s.t. $\forall \, y \in Y$, $\overline{ \overline{y}} = y$, $\overline{y} \neq y$, $o(y) = t(\overline{y})$.  

vertex $P \in X$ of $\Gamma$. 

(oriented) edge $y\in Y$, $\overline{y} \equiv $ inverse edge.  

origin of $y := $ vertex $o(y) = t(\overline{y})$.  

terminus of $y:= $ vertex $t(y) = o(\overline{y})$   

extremities of $y:= \lbrace o(y),t(y)\rbrace$  

If 2 vertices \textbf{adjacent}, they're extremities of some edge.  

orientation of graph $\Gamma = Y_+ \subset Y = \text{ edge } \Gamma$ s.t. $Y = Y_+ \coprod \overline{Y}_+$.  It always exists.  

oriented graph defined, up to isomorphism, by giving 2 sets $X,Y_+$ and $Y+ \to X\times X$.  
	
	corresponding set of edges is $Y = Y_+\coprod \overline{Y}_+$ where $\overline{Y}_+ \equiv $ copy of $Y_+$  
\end{definition}

\subsubsection{Realization of a Graph}

cf. Realization of a Graph in Serre (1980) \cite{agat:Serr1980}.  

Let graph $\Gamma$, $X = \text{vert}\Gamma$, $Y = \text{edge}\Gamma$.

topological space $T = X \coprod Y \times [0,1]$, where $X,Y$ provided with discrete topology.  

Let $R$ be finest equivalence relation on $T$ for which 
\begin{equation}
\begin{gathered} \begin{aligned}
	& (y,t) \equiv (\overline{y}, 1-t) \\ 
	& (y,0) \equiv o(y) \\
	& (y,1) \equiv t(y)
\end{aligned} \qquad \, \forall \, y \in Y, \, \forall \, t \in [0,1]
\end{gathered}
\end{equation}

quotient space $\text{real}(\Gamma) = T/R $ is \emph{realization} of graph $\Gamma$.  (realization is a functor which commutes with direct limits).  

Let $n\in \mathbb{Z}^+$.  Consider oriented graph of $n+1$ vertices $0,1,\dots n$,  \\

\begin{definition}
	path (of length $n$) in graph $\Gamma$ is morphism $c$ of $\mathbf{\text{Path}}_n$ into $\Gamma$  
\end{definition}

orientation given by $n$ edges $[i,i+1]$, $0\leq i <n$, $\begin{aligned} & \quad \\
& o([i,i+1]) = i \\
& t([i,i+1]) = i+1 \end{aligned}$  

For $n\geq 1$, \\
$(y_1\dots y_n)$ sequence of edges $y_i = c([i-1,i])$ s.t. 
\[
t(y_i) = o(y_{i+1}), \qquad \, 1 \leq i < n \text{ determine } c
\]
If $P_i = c(i)$,  \\
$c$ is a path from $P_0$ to $P_n$, and $P_0$ and $P_n$ are \emph{extremities of the path $c$}.  

pair of form $(y_i,y_{i+1})=(y_i, \overline{y}_i)$ in path is \textbf{backtracking}.  

path (of length $n-2$), from $P_0$ to $P_n$ given (for $n>2$) by $(y_1\dots y_{i-1}, y_{i+2}\dots y_n)$  

If $\exists \, $ path from $P$ to $Q$ in $\Gamma$, $\exists \, $ one without backtracking (by induction)  

direct limit $\mathbf{\text{Path}}_{\infty} = \varinjlim \mathbf{\text{Path}}_n$ provides notion of infinite path.  \\
$\mathbf{\text{Path}}_{\infty} \ni $ infinite sequence $(y_1,y_2 , \dots)$ of edges s.t. $t(y_i) = o(y_{i+1})$ \, $\forall \, i \geq 1$.  


\begin{definition}[connected graph; Def. 3 of Serre (1980) \cite{agat:Serr1980}]
	graph connected if $\forall \, $ 2 vertices, 2 vertices are extremities of at least 1 path.  
	
	maximal connected subgraphs (under relation of inclusion) are \emph{connected components} of graph.  
\end{definition}

\subsubsection{Circuits}  

Let $n\in \mathbb{Z}^+$, $n\geq 1$.  

Consider \\
set of vertices $\mathbb{Z}/n\mathbb{Z}$, orientation given by $n$ edges $[i,i+1]$, $(i\in \mathbb{Z}/n\mathbb{Z})$ with $\begin{aligned} & \quad \\
 & o([i,i+1]) = i \\
 & t([i,i+1]) = i+1 \end{aligned}$

\begin{definition}[circuit; Def. 4 of Serre (1980) \cite{agat:Serr1980}]
	circuit (length $n$) in graph is subgraph isormorphic to $\mathbf{\text{Circ}}_n$.  
\end{definition}
i.e. subgraph = path $(y_1\dots y_n)$, without backtracking, s.t. $P_i = t(y_i)$, \, $(1\leq i \leq n)$ distinct, s.t. $P_n = o(y_1)$

$n=1$ case: $\mathbf{\text{Circ}}_1$, $\mathbb{Z}/\mathbb{Z} = \lbrace 0 \rbrace$, $1$ edge, $[0,1]$, $0 \in \mathbb{Z}/1\mathbb{Z}$, $\begin{aligned} & \quad \\
	& o([0,1]) = 0 \\
	& t([0,1]) = 1 \end{aligned}$  
	
	Note $\mathbf{\text{Circ}}_1$ has automorphism of order 2, which changes its orientation, i.e. \\
	$\exists \, $ automorphism $\sigma \in \text{Aut}( \mathbf{\text{Circ}}_1) $ s.t. $|\sigma | = 2$, i.e. $\sigma^2=1$.  \\
	loop $:= $ circuit of length $1$; so loop $\in \overline{ \mathbf{\text{Circ}} }_1$.  
	
	path $(y_1)$, $P_1 = t(y_1) = o(y_1)$.  
	
	$n=2$ case: $\mathbf{\text{Circ}}_2$, $\mathbb{Z}/2\mathbb{Z} = \lbrace 0 ,1\rbrace$, 2 edges $[0,1], [1,2]$,  \\
	path $(y_1,y_2)$, $(1\leq i \leq 2)$, $\begin{aligned} & \quad \\
		& P_1 = t(y_1) \\
		& P_2 = t(y_2) = o(y_1) \end{aligned}$  
		
		
	

\subsection{Combinatorial graphs}

Let $(X,S)\equiv $ simplicial complex of dim. $\leq 1$, with \\
$X \equiv $ set \\
$S \equiv $ set of subsets of $X$ with $1$ or $2$ elements, containing all the 1-element subsets.  

associates with it a graph $\Gamma = (X, \lbrace (P,Q) \rbrace)$.  

$X$ is its set of vertices.  

edges $=\lbrace (P,Q)\in X\times X\rbrace$ s.t. $P\neq Q$, $\lbrace P ,Q \rbrace \in S$, with $\overline{(P,Q)} = (Q,P)$
\[
\begin{aligned}
& o(P,Q) = P \\
& t(P,Q) = Q
\end{aligned}
\]

In this graph, 2 edges with same origin and same terminus are equal.  This is equivalent to (see following Def.)

\begin{definition}[combinatorial; Def. 5 of Serre (1980) \cite{agat:Serr1980}]
	graph is combinatorial if it has no circuit of length $\leq 2$
\end{definition}
Conversely, it's easy to see that 

every combinatorial graph $\Gamma$ derived (up to isomorphism) by construction above from simplicial complex $(X,S)$, where \\
$X = \text{vert} \Gamma$ \\
$S=$ set of subset $\lbrace P,Q \rbrace$ of $X$ s.t. $P$ and $Q$ either adjacent or equal. 

